\NativeTitle{89}{总复习（三）}{（统计与概率）}{40}{4}{一 & 二 & 三 & 四}{ &  &  & }
\NativeSlot{282}{36}{%
\ExamGroup{一、填空。{\normalfont\kaishu\fontsize{10.8bp}{14bp}\selectfont （每空3分，共24分）}}

% p089_s01

\ExamQuestion{1.}{5}%
\mbox{（\hspace{1.8em}）}能较好地反映一组数据的整体水平。\par
}
\NativeSlot{332}{36}{%
% p089_s02

\ExamQuestion{2.}{5}%
小玉一分钟跳绳\(160\)个，小涵一分钟跳绳\(150\)个，他俩的平均跳绳个数是\mbox{（\hspace{1.8em}）}个，小涵再多跳\mbox{（\hspace{1.8em}）}个，他俩的跳绳数就一样多了。\par
}
\NativeSlot{385}{36}{%
% p089_s03

\ExamQuestion{3.}{5}%
四个数\(12,8,7,13\)的平均数是\mbox{（\hspace{1.8em}）}。\par
}
\NativeSlot{410}{36}{%
% p089_s04

\ExamQuestion{4.}{5}%
四年级\(3\)个班向希望小学捐赠图书，平均每班捐\(20\)本，其中，一班捐了\(23\)本，二班捐了\(15\)本，三班捐了\mbox{（\hspace{1.8em}）}本。\par
}
\NativeSlot{462}{36}{%
% p089_s05

\ExamQuestion{5.}{5}%
三个连续自然数的平均数是\(12\)，其中最大的数是\mbox{（\hspace{1.8em}）}。\par
}
\NativeSlot{488}{36}{%
% p089_s06

\ExamQuestion{6.}{5}%
在仰卧起坐达标测试中，\allowbreak{}\(A\)组\(3\)名同学共做\(96\)个，\allowbreak{}\(B\)组\(4\)名同学共做\(156\)个。平均每人做\mbox{（\hspace{1.8em}）}个。\par
}
\NativeSlot{539}{44}{%
% p089_s07

\ExamQuestion{7.}{5}%
绘画比赛有\(7\)个评委，分别给张梦打分\(98\)分、\(95\)分、\(96\)分、\(92\)分、\(95\)分、\(96\)分、\(99\)分，按照比赛规则，要去掉一个最高分和一个最低分，其他分数的平均数就是最后得分。张梦的最后得分是\mbox{（\hspace{1.8em}）}分。\par
}
\NativeSlot{617}{61}{%
\ExamGroup{二、选择（将正确答案的序号填在括号里）。{\normalfont\kaishu\fontsize{10.8bp}{14bp}\selectfont （每题2分，共10分）}}

% p089_s08

\ExamQuestion{1.}{5}%
一个人工湖测得最深处是\(2.7\)米，最浅处是\(1.1\)米，这个人工湖的平均水深可能是\mbox{（\hspace{1.8em}）}。\par
\ExamOptions{A. 2.7米}{B. 1.1米}{C. 1.9米}{D. 0.8米}
}
\NativeSlot{720}{42}{%
% p089_s09

\ExamQuestion{2.}{5}%
四年级一班学生的平均体重是\(35\)千克，小萱在四年级一班，她的体重\mbox{（\hspace{1.8em}）}是\(32\)千克。\par
\ExamOptions{A. 可能}{B. 不可能}{C. 一定}{D. 无法确定}
}
